\scanpage{20}
\mspage{5}
distance from both, these three half-circles joining at the two poles $P^+$ and $P^-$. The ``disc'' $Q_i$ around $s_i$ has a conical structure around $s_i$ (as has any conformal pointed disc), and we may take the concentric \deleted{\unsure{discs}} discs $\lambda_iQ_i$ with
\deleted{$\unsure{i}$}
\[
0<\lambda_i<1\ .
\]
The model I had in mind was that the (non degenerate) admissible systems of discs around the points $s_i$ ($i\in\{0,1,\infty\}$) \unsure{were} exactly the systems of discs $\lambda_iQ_i$, with $\lambda_i$ as above. (If we allow some discs to be degenerate, this means that instead of the unequality above we merely demand $0\le\lambda_i<1$, $0$ not excluded.)

This model, if correct, would give a rather precise description of the inclusion relationships between pants, when these are considered as embedded in the sphere. The intersection of all would be this system of three half circles $C_i$, and the two poles $P^+,P^-$ would play a significant role in the \unsure{geometry} of the pants, from this point of view. But it does'nt seem that neither those half circles (which need not be geodesical I guess), nor the two poles, have ever been described as intrinsically associated to a pant. Of course, this model would give alternative ``parameters'' $\lambda_i$ for describing a pant, which are best suited for grasping the pants in terms of spherical geometry. The next question would be an understanding of the relationship between these parameters, and Thurston's $\ell_i$. Maybe it is unreasonable to expect that for given index $i\in\{0,1,\infty\}$, the length $\ell_i$ depends only on $\lambda_i$ and not on the other parameters $\lambda_j$ - and for this reason, \deleted{\unsure{the}} the intuition at the beginning of this letter, using Thurston's coordinate functions and notably the $\ell_i$'s to get a fibration structure on $\widetilde{MB}$ over $\widetilde M$, in terms of a given decomposition of $X_o$ into pants, is probably not really relevant, namely it is non intrinsic. Assuming the model I am suggesting is correct, the accurate description \deleted{\unsure{but}} of $\widetilde{MB}$ in terms of $\widetilde M$ would be
\[
\widetilde{MB}\simeq\widetilde M\mathbin{\mathrm{x}}[0,1[^I\ ,
\]
where the second factor on the right hand sight refers to the system of multipliers $\lambda_i$ ($i\in I$), tied to the $r_i$ above by $r_i=\lambda_i\rho_i$.

My question of course is whether you have any information or idea to propose, especially on the basic problem of relying pants to spherical geometry, and more specifically, whether the model above is likely to hold, or is definitely false. Also, one difficulty we found with hyperbolic geometry of conformal surfaces,\inserted{is} that apart from existence and unicity of the hyperbolic structure (compatible with the given conformal one and for which the boundary is geodesic), there seems to be little
% archive foliation: 19
